\chapter{Matrice de traçabilité des exigences}
% Modèle du canevas E221, section 7. Une ligne par exigence (US-xx, BR-xx) du chapitre 2.
% Exigence -> composant -> test -> preuve -> statut. Une exigence sans test est marquée « non testée ».
% Sources : docs/requirements.md, docs/requirements/FEAT-*.md, docs/test-report.md, docs/evidence/.
% Pour une matrice longue, passer au paquet longtable.
% NB : les cellules commençant par « [ » sont entourées d'accolades (sinon \midrule lit un argument optionnel).

\begin{table}[htbp]
  \centering
  \caption{Matrice de traçabilité des exigences (modèle)}
  \label{tab:matrice-tracabilite}
  \begin{tabular}{@{}p{0.30\linewidth}p{0.30\linewidth}p{0.32\linewidth}@{}}
    \toprule
    \textbf{Exigence} & \textbf{Réalisation associée} & \textbf{Test et résultat} \\
    \midrule
    {[Identifiant et critère mesurable]} & {[Composant, fonctionnalité ou configuration]} & {[Identifiant du test, preuve, statut et limite]} \\
    \bottomrule
  \end{tabular}
\end{table}
